\documentclass[a4paper,12pt]{article}
\usepackage[utf8]{inputenc}
\usepackage[T1]{fontenc}
\usepackage[english]{babel}
\usepackage{geometry}
\geometry{left=2.5cm,right=2.5cm,top=2.5cm,bottom=2.5cm}

\usepackage{lmodern}
\usepackage{xcolor}
\definecolor{titleblue}{RGB}{0,51,140}

\usepackage{hyperref}

\pagestyle{empty}

\begin{document}

\begin{flushleft}
    \textbf{Ismail AISSAOUI} \\
    State Engineer in Artificial Intelligence \\
    ismail.aissaoui.pro@gmail.com \\
    +213 660 70 77 96 \\
    \href{https://ismail-aissaoui.vercel.app}{ismail-aissaoui.vercel.app}
\end{flushleft}

\begin{flushright}
    To the PhD Selection Committee \\
    \textbf{Uppsala University} \\
    Department of Information Technology \\
    \medskip
    April 30, 2026
\end{flushright}

\bigskip

\noindent \textbf{Subject: Application for PhD in Physics-Informed ML for Green Hydrogen Production}

\bigskip

Dear Members of the Selection Committee,

As a State Engineer in Artificial Intelligence from ESI-SBA, I am writing to express my enthusiastic application for the PhD position in physics-informed machine learning for green hydrogen production. My research focus on integrating physical laws into deep learning architectures for industrial systems makes me a perfect fit for this interdisciplinary project.

Currently, I am completing my thesis at Sonatrach, where I have developed a **Physics-Informed Digital Twin** for gas turbine monitoring. The core of my work involves the design of custom **PINN (Physics-Informed Neural Network)** loss functions that enforce thermodynamic consistency within generative surrogate models (Mamba-SSM/xLSTM). This experience has taught me how to handle the "noise vs. physics" trade-off in industrial sensors, a challenge that I am eager to tackle in the context of hydrogen production.

Furthermore, my expertise in **computerized image processing**—acquired through projects such as industrial safety monitoring (YOLOv8) and real-time vision systems (Recognizini)—provides me with the tools to implement the vision-based monitoring aspects of your project. I am particularly interested in how computerized image processing can provide real-time feedback for physics-based control loops in green energy production.

Uppsala University’s reputation for excellence in machine learning and control, combined with the collaboration with Alleima and Sandvik, offers the ideal environment for me to transition my skills toward sustainable energy. I am a highly autonomous researcher with a strong background in C++ and PyTorch, capable of delivering production-grade research prototypes.

Thank you for your time and consideration. I look forward to the possibility of discussing how my background in physics-aware AI can contribute to your research group.

Sincerely,

\bigskip

\textbf{Ismail AISSAOUI}

\end{document}
